\documentclass[11pt]{article}
\usepackage[margin=1in]{geometry}
\usepackage{amsmath,amssymb,amsthm,bm}
\usepackage{booktabs,array}
\usepackage[most]{tcolorbox}
\usepackage{listings}
\usepackage{tikz}
\usetikzlibrary{arrows.meta,positioning}
\usepackage{hyperref}
\usepackage{xcolor}
\hypersetup{colorlinks=true,linkcolor=blue!50!black,citecolor=blue!50!black,urlcolor=blue!50!black}

% ---- three-level boxes (same scheme as orbit_transfer/doc/algorithms_orbit_transfer.tex)
\newtcolorbox{eli}{colback=green!5,colframe=green!45!black,title=ELI5,fonttitle=\bfseries,breakable}
\newtcolorbox{intu}{colback=blue!4,colframe=blue!45!black,title=Intuition,fonttitle=\bfseries,breakable}
\newtcolorbox{rig}{colback=gray!6,colframe=black!70,title=Rigor,fonttitle=\bfseries,breakable}
\newtcolorbox{warn}{colback=orange!6,colframe=orange!70!black,title=Where it breaks / what is assumed,fonttitle=\bfseries,breakable}

\lstset{
  language=Matlab, basicstyle=\footnotesize\ttfamily, keywordstyle=\color{blue!60!black},
  commentstyle=\color{green!35!black}\itshape, stringstyle=\color{purple!60!black},
  numbers=left, numberstyle=\tiny\color{gray}, frame=single, rulecolor=\color{black!30},
  breaklines=true, showstringspaces=false, columns=fullflexible, xleftmargin=14pt
}

\newcommand{\file}[1]{\texttt{#1}}
\newcommand{\R}{\mathbb{R}}
\newcommand{\lam}{\lambda}
\newcommand{\tf}{t_f}

\title{Two Ways to Swing Up a Pendulum\\
\large Direct collocation and the indirect Pontryagin method, taught side by side}
\author{M.~Casey \quad (Coorbital, \file{collocation\_examples/})}
\date{2026-09-17}

\begin{document}
\maketitle

\begin{abstract}
The cart-pole swing-up in \file{collocation\_examples/} is solved twice in this
repository: once by \emph{direct collocation} (\file{ex2\_cart\_pole\_swing\_up/}),
which chops the trajectory into pieces and hands the whole thing to a nonlinear
programming solver, and once by the \emph{indirect} Pontryagin method
(\file{ex3\_cart\_pole\_pmp/}), which writes down the conditions an optimal
trajectory must satisfy and then solves those. The two land on the same motion:
their costs agree to $0.062\%$ and their control histories to $0.42\%$ RMS.

This document explains both methods, and the bridge between them --- the fact
that the multipliers the NLP solver produces \emph{are} samples of the costates
the Pontryagin method shoots. Every section is given three times: a picture with
no mathematics (ELI5), the reasoning an engineer needs in order to trust it
(Intuition), and the equations the code implements (Rigor). Code is quoted from
the repository as it actually stands, with line-level explanation.

A closing chapter is about a sign error that lived in this code for a long time,
what it cost, and the one test that would have caught it on day one.
\end{abstract}

\tableofcontents

% =======================================================================
\section{What problem is this, exactly?}

Before any method: \emph{which} optimal control problem is being solved? The
three classical choices lead to very different mathematics, and calling this one
by the wrong name will send you looking for structure that is not there.

\begin{eli}
We want to swing a pendulum from hanging straight down to balanced straight up,
in exactly five seconds, by sliding the cart it is attached to back and forth.
Lots of different pushes would get it there. We want the \emph{gentlest} one ---
the one that uses the least force overall. We are not in a hurry (five seconds
is fixed, not something we are trying to shrink), and we are not paying for fuel
by the litre. We are paying for \emph{effort}.
\end{eli}

\begin{intu}
Three costs, three worlds:
\begin{itemize}
\item \textbf{Minimum time} $J=\tf$: ``get there as fast as possible''. The final
  time is an unknown, and the answer almost always pushes as hard as the limits
  allow. The interesting structure is \emph{when} to switch between limits.
\item \textbf{Minimum fuel} $J=\int|u|\,dt$: ``spend as little propellant as
  possible''. Because $|u|$ is linear in the control's magnitude, the optimum
  is bang--bang: full on or full off, with a switching function deciding which.
  This is the world where you need saltation matrices, smoothing homotopies, and
  all the switch-detection machinery in \file{orbit\_transfer/}.
\item \textbf{Minimum energy (effort)} $J=\int u^2\,dt$: ``push as gently as
  possible''. The cost is \emph{quadratic}, so the optimal control comes out as a
  smooth formula rather than a switching rule. Nothing bangs; nothing switches.
\end{itemize}
Ours is the third, with $\tf$ fixed at $5$~s. That choice is what makes this
problem a good teaching example: the optimality condition can be solved for $u$
in closed form, so the only hard part left is the boundary-value problem ---
which is exactly the part we want to teach.
\end{intu}

\begin{rig}
The problem solved by both \file{ex2} and \file{ex3} is
\[
\min_{u(\cdot)}\; J=\int_0^{\tf} u(t)^2\,dt,
\qquad \tf = 5~\text{s fixed},
\]
subject to $\dot x = f(x)+g(x)\,u$ with
$x=[q_1,\;q_2,\;\dot q_1,\;\dot q_2]^\top$ and boundary conditions
\[
x(0)=[0,\;0,\;0,\;0]^\top,\qquad x(\tf)=[0,\;\pi,\;0,\;0]^\top ,
\]
i.e.\ hanging at rest to upright at rest, cart returning to the origin. The
control is unconstrained: no bound enters the formulation. (The direct fixture
is solved at a deliberately loose bound of $2000$~N purely so that the bound is
inactive --- measured $\max|u| = 68.5$~N --- and it therefore approximates the
same unconstrained optimum.)

Naming: in the literature this is \emph{minimum energy} or \emph{minimum
effort}; both names appear. ``Energy'' here refers to the \emph{control} energy
$\int u^2$, which is not the mechanical energy of the cart-pole. The mechanical
energy makes its own appearance in \S\ref{sec:signerror}, as a test.
\end{rig}

\begin{warn}
The word ``energy'' is overloaded. $\int u^2\,dt$ has units of
$\mathrm{N^2\,s}$ --- it is not joules, and it is not the work done by the
force ($\int u\,\dot q_1\,dt$, which \emph{is} joules and \emph{is} the quantity
the physics test in \S\ref{sec:signerror} uses). Keep them apart.
\end{warn}

% =======================================================================
\section{The plant: a cart, a rod, a bob}

\begin{eli}
A cart slides on a rail. A rigid rod hangs off it with a weight on the end. You
can push the cart left or right; you cannot touch the rod. Everything the rod
does, it does because the cart moved underneath it --- like balancing a broom on
your palm, except your palm may only slide along a straight line.
\end{eli}

\begin{intu}
Two degrees of freedom ($q_1$, the cart's position; $q_2$, the rod's angle) and
only one actuator (the force $u$ on the cart). That deficit --- more things to
control than knobs to control them with --- is what makes this
\emph{underactuated}, and it is why the answer has to be a swing: the only way to
put energy into the rod is to move the cart so that the rod is dragged, and the
only way to stop the rod at the top is to move the cart out from under it at
exactly the right moment.

The single most important structural fact for everything that follows is that
the force enters \emph{linearly}. Push twice as hard and the acceleration
contribution doubles exactly. Formally $\dot x = f(x)+g(x)u$: a part that happens
anyway (the drift $f$, gravity and centrifugal terms) plus a part proportional
to what you do ($g$, the control column). Almost every nice property we exploit
downstream traces back to that split.
\end{intu}

\begin{rig}
With $q_2$ measured from \emph{hanging}, the bob sits at
$(q_1+L\sin q_2,\;-L\cos q_2)$, and the Lagrangian equations give
\[
f(x)=\begin{bmatrix}\dot q_1\\[2pt] \dot q_2\\[4pt]
\dfrac{Lm_2 s\,\dot q_2^{\,2}+m_2 g\,c\,s}{D_1}\\[10pt]
-\dfrac{Lm_2 c\,s\,\dot q_2^{\,2}+(m_1+m_2)g\,s}{D_2}\end{bmatrix},
\qquad
g(x)=\begin{bmatrix}0\\0\\ \dfrac{1}{D_1}\\[8pt] -\dfrac{c}{D_2}\end{bmatrix},
\]
with $s=\sin q_2$, $c=\cos q_2$, $D_1=m_1+m_2(1-c^2)=m_1+m_2 s^2$ and
$D_2=L(m_1+m_2)\bigl(1-\tfrac{m_2}{m_1+m_2}c^2\bigr)=L\,D_1$. Constants:
$m_1=5$~kg, $m_2=1$~kg, $L=2$~m, $g=9.8~\mathrm{m/s^2}$.

Note the minus signs on the fourth row of both $f$ and $g$. They are the subject
of \S\ref{sec:signerror}.

In code, \file{ex3\_cart\_pole\_pmp/cartpole\_field.m} is exactly this split:

\begin{lstlisting}
q1d = x(3);  q2d = x(4);
s = sin(x(2));  c = cos(x(2));
D1 = p.m1 + p.m2*(1 - c^2);
D2 = p.L*(p.m1 + p.m2)*(1 - (p.m2/(p.m1 + p.m2))*c^2);

F = [q1d;
     q2d;
     (p.L*p.m2*s*q2d^2 + p.m2*p.g*c*s)/D1;
    -(p.L*p.m2*c*s*q2d^2 + (p.m1 + p.m2)*p.g*s)/D2];
G = [0; 0; 1/D1; -c/D2];
\end{lstlisting}

Two details that matter later. First, \texttt{F} is the derivative at $u=0$ and
\texttt{G} is exactly $\partial\dot x/\partial u$, so \texttt{F + G*u} is the full
dynamics --- the test suite checks this against an independently written copy of
the equations rather than trusting the claim. Second, the file contains no
\texttt{abs}, \texttt{norm}, \texttt{max}, \texttt{min}, or branch on a state
value. That makes it \emph{complex-step safe}, which \S\ref{sec:stm} cashes in.
\end{rig}

% =======================================================================
\section{Method 1: direct collocation (discretize, then optimize)}

\begin{eli}
Guess the whole movie at once. Chop the five seconds into 200 slices and write
down, for every slice boundary, where the cart is, how fast it is going, where
the rod is pointing, and how hard you are pushing. That is a big list of
numbers --- about a thousand of them. Now hand the list to a solver and say:
``adjust these until (a) the pieces fit together like a physically possible
movie, (b) the start and end look the way I asked, and (c) the total push is as
small as you can make it.''

The solver does not know any physics. It only knows the rules you wrote down.
\end{eli}

\begin{intu}
The trick that makes this work is how ``physically possible'' is expressed. You
do \emph{not} integrate the equations of motion. Instead you write, for each
slice, an equation that says: \emph{the change in the state across this slice
equals the average of the slopes at its two ends, times the slice width}. That
is the trapezoid rule, turned from an integration scheme into a constraint.

Each such equation is called a \emph{defect}: it is zero when the piece is
consistent, nonzero when it is not. The solver's job is to drive all of them to
zero while minimizing cost. The states and controls are all decision variables
on an equal footing, which is why this family is called
\emph{simultaneous} or \emph{discretize-then-optimize}.

Why it is popular: it is robust, it needs no good initial guess for the
costates (there are none), and bounds on states or controls are just bounds on
variables. Why it has a ceiling: the answer is only as good as the mesh, and the
optimality it achieves is the optimality of the \emph{discretized} problem,
which is not quite the problem you asked about.
\end{intu}

\begin{rig}
Let $N$ be the number of segments, $h=\tf/N$, and let $x_k, u_k$ denote the
state and control at node $k=1,\dots,N+1$. The decision vector stacks them by
variable:
\[
Z=\bigl[q_1^{(1..N+1)},\;q_2^{(1..N+1)},\;\dot q_1^{(1..N+1)},\;\dot q_2^{(1..N+1)},\;u^{(1..N+1)}\bigr]^\top\in\R^{5(N+1)} .
\]
The trapezoidal defect for segment $k$ is
\[
\zeta_k \;=\; x_{k+1}-x_k-\frac{h}{2}\bigl[\,\hat f(x_{k+1},u_{k+1})+\hat f(x_k,u_k)\,\bigr] \;=\;0,
\qquad \hat f = f+g\,u,
\]
and the cost is the trapezoid quadrature of the running cost,
\[
J\approx\frac{h}{2}\sum_{k=1}^{N}\bigl(u_k^2+u_{k+1}^2\bigr).
\]
Together with the eight boundary conditions this is a nonlinear program in
$\R^{5(N+1)}$ with $4N+8$ equality constraints, handed to \texttt{fmincon}.

The defects, from \file{ex2\_cart\_pole\_swing\_up/try2/nonlcon\_cart\_pole.m},
written vectorized over all segments at once:

\begin{lstlisting}
% kinematic rows: change in position = average velocity x h
ceq1 = q1(2:end) - q1(1:end-1) - h/2.*(q1dot(2:end) + q1dot(1:end-1));
ceq2 = q2(2:end) - q2(1:end-1) - h/2.*(q2dot(2:end) + q2dot(1:end-1));

% dynamic rows: change in velocity = average acceleration x h
term_kp1_cart_accel = cart_accel(q2(2:end), q2dot(2:end), u(2:end), L, m1, m2, g);
term_k_cart_accel   = cart_accel(q2(1:end-1), q2dot(1:end-1), u(1:end-1), L, m1, m2, g);
ceq3 = q1dot(2:end) - q1dot(1:end-1) - h/2.*(term_kp1_cart_accel + term_k_cart_accel);
\end{lstlisting}

Line by line: \texttt{q1(2:end)} is every node from the second onward and
\texttt{q1(1:end-1)} every node up to the second-to-last, so their difference is
the vector of changes across all $N$ segments in one shot. Subtracting
\texttt{h/2} times the sum of the endpoint slopes is precisely $\zeta_k$ for
every $k$. Rows \texttt{ceq1}/\texttt{ceq2} use velocity as the slope of
position (that part is exact kinematics); \texttt{ceq3}/\texttt{ceq4} call the
acceleration helpers for the slope of velocity.

And the objective, from \file{objective.m}:

\begin{lstlisting}
J = h/2 .* sum (u(2:end).^2 + u(1:end-1).^2,1);
\end{lstlisting}

which is the quadrature above, again vectorized over segments.
\end{rig}

\begin{warn}
A converged NLP means \emph{the defects are small}, not \emph{the trajectory is
accurate}. These are different claims, and on stiff problems they can differ by
many orders of magnitude --- the orbit-transfer library has a documented case of
a defect of $1.4\times10^{-14}$ next to a true local error of $1.5$. The honest
check is to re-integrate the reconstructed control with a real integrator and
see where you actually arrive; \file{oc.local\_residual} and
\file{oc.fly\_control} exist for exactly this.
\end{warn}

% =======================================================================
\section{Method 2: the indirect method (optimize, then discretize)}

\begin{eli}
Instead of guessing the whole movie, ask a different question: \emph{what must
be true of the best movie?} Physics and calculus answer that for you. It turns
out that alongside the four numbers describing the cart and rod, there are four
more shadow numbers --- one attached to each --- that tell you how much the final
result would improve if that quantity were nudged. Those shadows obey their own
equations, running alongside the real ones.

Here is the magic: once you know the four shadows at the \emph{start}, the best
push at every instant is given by a formula, and the whole movie follows by
playing the equations forward. So the entire problem shrinks from a thousand
unknowns to four. The catch is that you do not know those four numbers --- you
have to hunt for the ones that make the rod finish upright.
\end{eli}

\begin{intu}
The shadow variables are the \emph{costates} $\lambda$, and they are Lagrange
multipliers for the dynamics --- the price of the constraint ``the state must
actually obey the equations of motion''. Pontryagin's Minimum Principle says an
optimal trajectory makes a certain scalar, the Hamiltonian
\[
H = \underbrace{u^2}_{\text{what this instant costs}} \;+\; \underbrace{\lambda^\top(f+gu)}_{\text{what this instant buys}},
\]
as small as possible \emph{at each instant}, with respect to $u$ alone.

Because our cost is quadratic in $u$ and the dynamics are linear in $u$, $H$ is a
quadratic in $u$ with a positive leading coefficient --- a parabola opening
upward. Minimizing a parabola is done by setting its derivative to zero, which
gives $u$ in closed form. This is the sense in which minimum-energy problems are
\emph{easy}: there is no switching, no case analysis, no bang--bang. Change the
cost to $\int|u|$ and the parabola becomes a V with a flat bottom, and all the
switching machinery appears.

So the method is: eliminate $u$ analytically, leaving an eight-dimensional flow
in $(x,\lambda)$; then find the four unknown initial costates that make the
terminal conditions come out right. Four unknowns, four equations. That is a
root-find, not an optimization --- which is why it converges quadratically when
it converges, and why it can be delicate.
\end{intu}

\begin{rig}
With $H=u^2+\lambda^\top(f+gu)$, the three PMP conditions are
\[
\dot x = \frac{\partial H}{\partial\lambda}=f+gu,
\qquad
\dot\lambda = -\frac{\partial H}{\partial x},
\qquad
u^\star=\arg\min_u H .
\]
Stationarity gives $\partial H/\partial u = 2u+\lambda^\top g = 0$, hence
\[
\boxed{\,u^\star=-\tfrac12\lambda^\top g\,}
\qquad\text{and}\qquad
\frac{\partial^2H}{\partial u^2}=2>0,
\]
so the stationary point is the \emph{pointwise minimizer} everywhere (the
Legendre--Clebsch condition holds strictly). Substituting $u^\star$ closes an
autonomous 8-state system in $y=[x;\lambda]$.

The costate equation deserves care. After substituting $u^\star(x,\lambda)$ the
Hamiltonian becomes $\bar H(x,\lambda)=H(x,\lambda,u^\star)$, and
\[
\bar H_x = H_x + H_u\,u^\star_x = H_x
\qquad\text{because } H_u=0 \text{ at } u^\star .
\]
So one may differentiate \emph{holding $u$ fixed} and then evaluate at $u^\star$
--- the $\partial u^\star/\partial x$ term drops out exactly. Since the running
cost $u^2$ has no state dependence,
\[
\dot\lambda = -A(x,u^\star)^\top\lambda,
\qquad A = \frac{\partial (f+gu)}{\partial x}\bigg|_{u\ \text{fixed}} .
\]
Boundary conditions: $x(0)$ is fully specified (4 conditions) and $x(\tf)$ is
fully specified (4 more). With $\tf$ fixed there is no transversality condition
on $\tf$, and because the terminal state is fully constrained there is no
condition $\lambda(\tf)=0$ either. The unknowns are exactly $\lambda(0)\in\R^4$.

The whole of \file{cartpole\_pmp\_rhs.m} is four lines:

\begin{lstlisting}
x = y(1:4);  lam = y(5:8);
[F, G] = cartpole_field(x, p);
u = -(lam.'*G)/2;
A = cartpole_state_jac(x, u, p);
dy = [F + G*u; -A.'*lam];
\end{lstlisting}

Line 3 is $u^\star$. Line 4 evaluates $A$ at that control --- fixed-$u$
differentiation, per the argument above. Line 5 stacks $\dot x$ on top of
$\dot\lambda$. Note \texttt{.'} (transpose) rather than \texttt{'} (conjugate
transpose) in both places: on a complex-stepped path a conjugate transpose would
silently destroy the derivative.
\end{rig}

\begin{warn}
$H_{uu}>0$ makes $u^\star$ the unique minimizer \emph{at each instant}. It does
\emph{not} make the resulting trajectory a minimizing one. The boundary-value
problem can have several extremals, saddles among them; ruling those out needs a
second-order test (a conjugate-point or Jacobi test --- see
\file{costate\_common/ms\_conjugate\_test.m} in the orbit-transfer library).
Everything this demo verifies is consistent with a \emph{normal extremal},
which is a weaker statement than ``the optimum''.
\end{warn}

% =======================================================================
\section{The bridge: NLP multipliers \emph{are} costates}\label{sec:bridge}

This is the conceptual heart of the pairing, and the reason the two examples
live next to each other.

\begin{eli}
The first method never mentions shadow numbers. But it computes them anyway,
without being asked --- they fall out of how the solver enforces ``the pieces
must fit''. So you can run the easy, robust method, read the shadows out of its
answer, and hand them to the fussy, precise method as a starting guess. Easy
method warms up the hard one.
\end{eli}

\begin{intu}
When \texttt{fmincon} enforces the defect $\zeta_k=0$, it attaches a Lagrange
multiplier $\mu_k$ to it. Stationarity of the NLP's own Lagrangian with respect
to the state variables is, term by term, a discretization of the adjoint
equation $\dot\lambda=-H_x$. So $\mu_k$ is a sample of $\lambda$ --- up to a
scaling that depends on how the defect was written, a sign that depends on the
solver's conventions, and a \emph{station}: for the trapezoid rule the natural
association is with nodes, for Hermite--Simpson with interval midpoints.

Getting that association wrong injects a half-step phase error into every seed,
and it is a genuinely easy mistake --- it has been made in this codebase and
caught by review. That is why the rules live in one place,
\file{oclib/+oc/duals\_to\_costates.m}, rather than being re-derived per campaign.
\end{intu}

\begin{rig}
The seed construction in \file{run\_cartpole\_pmp.m}:

\begin{lstlisting}
[lamS, tS] = oc.duals_to_costates(struct('scheme', 'trapezoid', 'mu', R.muDefect, ...
                                         'tNodes', R.tN, 'velRows', 3:4));
\end{lstlisting}

\texttt{'scheme','trapezoid'} selects the station rule; \texttt{mu} is the
$4\times N$ array of defect multipliers harvested from the direct solve;
\texttt{velRows} names which rows hold the velocity costates (the default
\texttt{4:6} is for the 7-state CR3BP problems, and would be wrong here).

The global sign is convention-dependent. The library resolves it by a ``primer
vote'': the implied thrust direction $-\lambda_v/\|\lambda_v\|$ should align
with the thrust direction the solver actually used. A scalar force has no
direction to vote on, so this demo does the scalar analogue --- and, after
review, does it with a confidence measure rather than a bare sign test:

\begin{lstlisting}
signCorr = sum(uImplied.*uDirect) / sqrt(sum(uImplied.^2)*sum(uDirect.^2));
assert(abs(signCorr) > 0.5, 'run_cartpole_pmp:sign', ...
       'implied vs direct control correlation %.3f: the seed carries no usable sign information', signCorr);
if signCorr < 0, lamS = -lamS;  uImplied = -uImplied; end
\end{lstlisting}

The normalized correlation is a real diagnostic: measured $0.9998$ here. A bare
$\sum u_{\text{implied}}u_{\text{direct}}\neq 0$ test would have accepted an
arbitrarily weak agreement. The code also reports the \emph{amplitude ratio}
($1.009$ measured), because a sign vote cannot detect a costate seed that points
the right way but is scaled wrongly.
\end{rig}

% =======================================================================
\section{Multiple shooting and the state-transition matrix}\label{sec:stm}

\begin{eli}
Playing the equations forward for five seconds from a guessed start is like
rolling a marble down a long bumpy hill and hoping it lands in a cup: tiny
errors at the top become huge misses at the bottom. So we cut the five seconds
into eight shorter rolls, let each one start wherever it likes, and add the rule
``where one roll ends is where the next begins''. More unknowns, but each roll is
short enough to be predictable.

To steer the hunt we also need to know how the landing point moves when the
starting point moves. That sensitivity is a matrix, and we compute it by
carrying it along with the trajectory.
\end{eli}

\begin{intu}
The reason single shooting fails on long arcs is Lyapunov amplification: the
sensitivity of the endpoint to the initial costate grows roughly exponentially
with the arc length. In the orbit campaigns, catalog seeds that are perfectly
good miss by 36{,}000 to 560{,}000~km when single-shot. Splitting the horizon
into $K$ segments replaces one badly-conditioned map with $K$ well-conditioned
ones, at the price of carrying the junction states as extra unknowns and adding
continuity constraints.

The Newton step needs the Jacobian of ``miss'' with respect to the unknowns, and
the pieces of that Jacobian are \emph{state-transition matrices}: $\Phi =
\partial y(t_{k+1})/\partial y(t_k)$. You get $\Phi$ by integrating its own
differential equation alongside the state, $\dot\Phi = A_8\Phi$, starting from
the identity.
\end{intu}

\begin{rig}
The engine (\file{oclib/+oc/ms\_bvp.m}) asks the problem for three closures:
\texttt{prop} (propagate, optionally with $\Phi$), \texttt{rhs} (the field at a
point), and \texttt{terminal} (the terminal residual and its Jacobian). The demo
supplies them:

\begin{lstlisting}
prob = struct('ny', 8, 'freeIdx0', 5:8, ...
    'prop', @(dt, y0, needSTM) cartpole_pmp_prop(dt, y0, needSTM, p), ...
    'rhs',  @(y) cartpole_pmp_rhs(y, p), ...
    'terminal', @(y, needJ) terminalFcn(y, xf, needJ));
\end{lstlisting}

\texttt{'ny',8} is the augmented state dimension; \texttt{'freeIdx0',5:8} says
that of the eight components at $t=0$, only rows 5--8 (the costates) are
unknowns --- rows 1--4 are the fixed initial state. With \texttt{opts.fixedTf}
the unknown vector is
$p=[\lambda(0);\,y_2;\dots;y_K]$ and the equations are the $K-1$ continuity
blocks plus the 4 terminal conditions.

The variational integration, from \file{cartpole\_pmp\_prop.m}:

\begin{lstlisting}
y = z(1:8);
PHI = reshape(z(9:72), 8, 8);
dy = cartpole_pmp_rhs(y, p);
A = zeros(8);
h = 1e-20;
for col = 1:8
    yp = y;  yp(col) = yp(col) + 1i*h;
    A(:,col) = imag(cartpole_pmp_rhs(yp, p))/h;
end
dz = [dy; reshape(A*PHI, 64, 1)];
\end{lstlisting}

Seventy-two equations: eight for $y$, sixty-four for $\Phi$. The Jacobian $A$
comes from a \emph{complex step}: perturb one component by $ih$ with
$h=10^{-20}$ and read the imaginary part of the output. Unlike a finite
difference this has no subtractive cancellation, so it is accurate to machine
precision even at absurdly small $h$.

Which raises the subtlety that dictated the design of
\file{cartpole\_state\_jac.m}. The complex step above runs \emph{through}
\file{cartpole\_pmp\_rhs}, which itself needs $A$ for the costate equation. Had
that inner $A$ also been computed by a complex step, the two would share the
same imaginary channel and the inner derivative would be silently corrupted ---
no error, no warning, just wrong numbers. That is why the inner Jacobian is
generated symbolically instead:

\begin{lstlisting}
syms q1 q2 q1d q2d u m1 m2 L g real
xs = [q1; q2; q1d; q2d];
ps = struct('m1', m1, 'm2', m2, 'L', L, 'g', g);
[F, G] = cartpole_field(xs, ps);        % the FIELD ITSELF, on symbols
A = simplify(jacobian(F + G*u, xs));
matlabFunction(A, 'File', out, 'Vars', {xs, u, [m1; m2; L; g]}, ...);
\end{lstlisting}

The generator calls \file{cartpole\_field} on symbolic inputs, so there is no
second copy of the physics to drift out of step --- a lesson learned the hard
way (\S\ref{sec:signerror}).
\end{rig}

% =======================================================================
\section{What the two methods produce}

Both solve the same problem; the comparison is the point.

\begin{center}
\begin{tabular}{lrr}
\toprule
& \textbf{Direct (ex2 formulation)} & \textbf{Indirect (ex3)}\\
\midrule
unknowns & $1005$ & $4$ (plus $7\times8$ junction states)\\
cost $J$ & $2781.109846$ & $2779.381719$\\
terminal miss & pinned by constraint & $3.67\times10^{-11}$\\
needs a costate guess & no & yes\\
handles bounds & trivially & needs switching structure\\
accuracy limit & the mesh & the integrator\\
\bottomrule
\end{tabular}
\end{center}

The costs differ by $0.062\%$, and the control histories by $0.42\%$ RMS. The
gap is the direct method's discretization error: the indirect solve is the more
accurate of the two, because it never discretizes the dynamics at all --- it
integrates them.

\begin{intu}
The practical division of labour, and the reason the orbit-transfer library uses
both: the direct method is \emph{robust but coarse}, the indirect method is
\emph{sharp but fragile}. So run the direct method to get into the right basin,
harvest its multipliers (\S\ref{sec:bridge}), and let the indirect method polish
the answer to integrator accuracy. That is precisely the pipeline the costate
catalogs use, at scale, on orbit transfers.
\end{intu}

% =======================================================================
\section{How we know it is right}

\begin{eli}
Getting a pretty answer is easy. Knowing it is the right answer is the job. We
check four different things, and each one can fail on its own.
\end{eli}

\begin{intu}
Verification has layers, and they catch different faults:
\begin{enumerate}
\item \textbf{Did the root-find converge?} The engine's own residual. Necessary,
  and nearly meaningless by itself.
\item \textbf{Does the trajectory hit the target?} Re-fly $\lambda(0)$ from
  scratch and measure the terminal miss. This catches a converged solve of the
  \emph{wrong} equations.
\item \textbf{Is it an extremal?} Along an autonomous, fixed-$\tf$ extremal the
  Hamiltonian is \emph{constant} (not zero --- that would require free $\tf$).
  Measured spread: $5.8\times10^{-13}$ relative.
\item \textbf{Does the control actually fly?} Take $u(t)$ and integrate the true
  dynamics with it. Measured arrival error: $1.7\times10^{-7}$.
\end{enumerate}
A trap worth naming: a check that recomputes the thing it is testing proves
nothing. The stationarity residual $|2u+\lambda^\top g|$ is exactly zero when $u$
was \emph{built} as $-\lambda^\top g/2$ --- it measures arithmetic, not
optimality. This demo used to report that number as evidence; it no longer does.
\end{intu}

\begin{rig}
The engine's own residual and the reported miss are different quantities and
were confused here once. Measured on the converged root: the engine's last-arc
terminal residual is $7.27\times10^{-14}$, while re-flying $\lambda(0)$ over the
whole horizon gives $6.88\times10^{-9}$ at \texttt{RelTol} $10^{-12}$,
$8.20\times10^{-10}$ at $10^{-13}$, and $3.67\times10^{-11}$ at $10^{-14}$. The
loose figure was the \emph{measurement}, not the solution --- the reporting
integration's own error. The fix was to tighten the measurement, not the gate.
\end{rig}

% =======================================================================
\section{A sign error, and the test that would have caught it}\label{sec:signerror}

\begin{eli}
For a long time the equations in this repository had one minus sign missing. It
made the hanging position \emph{unstable} --- in that world, a pendulum hanging
straight down would fall \emph{upwards}. Both methods solved that world happily,
and agreed with each other beautifully, because they were both asking the same
wrong question. Agreeing with yourself is not evidence.
\end{eli}

\begin{intu}
Every test in the folder was \emph{relative}: the field was checked against a
copy of the equations, the Jacobian against the field, the STM against the flow,
the indirect solve against the direct one. A fault shared by all of them cancels
everywhere and is invisible.

What breaks the circle is a test whose reference comes from \emph{outside} the
equations under test. For mechanics the cheapest such reference is a
conservation law. You can write the energy from the \emph{geometry} alone ---
where the masses are and how fast they move --- without consulting the
acceleration formulas at all. Then physics demands a relationship between that
energy and the work being done, and any sign error in the accelerations breaks
it.
\end{intu}

\begin{rig}
Build the energy from positions only. The bob is at
$(q_1+L\sin q_2,\,-L\cos q_2)$, so
\[
E=\tfrac12 m_1\dot q_1^2+\tfrac12 m_2\bigl\|\dot r_{\text{bob}}\bigr\|^2-m_2gL\cos q_2 ,
\]
and the only external force doing work is $u$ acting on the cart, giving the
\emph{pointwise} identity
\[
\frac{dE}{dt}=\nabla E\cdot \dot x = u\,\dot q_1 \qquad\text{for every } (x,u).
\]
No integration is needed: evaluate both sides at a scatter of states and
controls. From \file{tests/test\_cartpole\_physics.m}:

\begin{lstlisting}
gradE = energy_gradient(x, p);          % complex step through the GEOMETRIC energy
[F, G] = cartpole_field(x, p);
xdot = F + G*u;
worstField = max(worstField, abs(gradE.'*xdot - u*x(3)));
\end{lstlisting}

Measured: $5.7\times10^{-14}$~W against a $460$~W scale with the corrected
equations, and $337$~W with the old ones. The same file also checks that the
linearization about $q_2=0$ is oscillatory with
$\omega^2=(m_1+m_2)g/(Lm_1)=5.880000$ (hanging is \emph{stable}), that $q_2=\pi$
is a real unstable pair, and that an unforced flight conserves $E$
($1.6\times10^{-12}$~J over 3~s).
\end{rig}

\begin{warn}
The general rule, worth carrying to any physics code: \textbf{at least one test
must not be a mirror}. If every check compares the implementation against
another expression of the same implementation, the suite measures internal
consistency and nothing else. Conservation laws, analytic limits, equilibrium
character and dimensional analysis are all cheap sources of an outside opinion.
\end{warn}

% =======================================================================
\section{Where the code lives}

\begin{center}
\begin{tabular}{ll}
\toprule
\textbf{path} & \textbf{what it is}\\
\midrule
\file{ex2\_cart\_pole\_swing\_up/try2/} & the direct method: objective, defects, dynamics helpers\\
\file{ex3\_cart\_pole\_pmp/cartpole\_field.m} & $f$ and $g$, the one home for the physics\\
\file{ex3\_cart\_pole\_pmp/cartpole\_pmp\_rhs.m} & the 8-state PMP field\\
\file{ex3\_cart\_pole\_pmp/cartpole\_pmp\_prop.m} & propagator and variational STM\\
\file{ex3\_cart\_pole\_pmp/run\_cartpole\_pmp.m} & seed, shoot, verify, report\\
\file{ex3\_cart\_pole\_pmp/run\_tests.m} & the runner that fails on failure\\
\file{oclib/+oc/ms\_bvp.m} & the shared multiple-shooting engine\\
\file{oclib/+oc/duals\_to\_costates.m} & the covector rules (\S\ref{sec:bridge})\\
\file{orbit\_transfer/doc/algorithms\_orbit\_transfer.tex} & the same ideas at campaign scale\\
\bottomrule
\end{tabular}
\end{center}

\paragraph{Reading order for someone new.} Run \file{run\_cartpole\_pmp} and
watch it print; read \file{cartpole\_pmp\_rhs.m} (four lines, the whole
principle); then \file{run\_cartpole\_pmp.m} top to bottom; then
\file{test\_cartpole\_physics.m}, which is the shortest statement of what it
means to actually believe a dynamics model.

% =======================================================================
\section{The executable companion, and the objective that never has to choose}\label{sec:companion}

\begin{eli}
Everything this document described is also a script you can run:
\file{cartpole\_minenergy\_study.m}. Press Run and it works through the same
story in the same order, printing PASS or FAIL as it goes instead of only at
the end. One more thing worth noticing on the way out: at no point did the
best push ever require a decision. It fell out of a formula, every time, at
every instant. The next two chapters of this project --- not written yet ---
change exactly that: give the cart a top speed of push it may not exceed, or
make it pay for every newton instead of every newton squared, and the same
question stops having a formula for an answer and starts having a switch.
\end{eli}

\begin{intu}
The script is not organized section-for-section with this document --- it is
organized as a pipeline, load $\to$ seed $\to$ solve $\to$ check --- so here is
the correspondence:
\begin{itemize}\itemsep2pt
\item \S1 (what problem is this) $\leftrightarrow$ the script's steps 0 and 2:
  the tolerance table, then the same cost and boundary conditions, in code.
\item \S2 (the plant) $\leftrightarrow$ step 1: \file{cartpole\_params}, and
  the physics oracle (V1) run and asserted \emph{first}, before anything
  downstream is allowed to mean anything.
\item \S4--\S\ref{sec:stm} (the indirect method, the bridge, shooting)
  $\leftrightarrow$ steps 3 and 4: the seed lifted verbatim from
  \file{run\_cartpole\_pmp.m}'s own multiplier block, then \texttt{oc.ms\_bvp}.
\item \S7 (what the two methods produce) $\leftrightarrow$ the X1 check
  folded into step 6.
\item \S8 (how we know it is right) $\leftrightarrow$ steps 5--8: the
  independent-seed check (X2), the necessary conditions (N1--N6), the
  sufficiency test (S1--S2), and the cross-check against the library
  instrument.
\item \S\ref{sec:signerror} (the sign error) $\leftrightarrow$ step 1's V1
  oracle again --- the same outside-the-equations check, run once and relied
  on everywhere after.
\end{itemize}
Why no switch here: an unconstrained control being minimized against a
\emph{strictly convex} cost is a parabola with nowhere to hit a wall. Put a
wall in front of it (a bound on $u$) or straighten the parabola into a line
or a V (a cost linear in $u$ or in $|u|$) and the minimizer stops living where
calculus finds it and starts living at the wall --- which is a decision, made
anew whenever the wall relocates.
\end{intu}

\begin{rig}
Dropping the parts of $H$ that do not depend on $u$, the three costs reduce
to $H(u) = u^2 + \sigma u$ here, $H(u) = 1 + \sigma u$ for minimum time, and
$H(u) = |u| + \sigma u$ for minimum fuel, all sharing $\sigma = \lambda^\top g$
--- the switching function already named in \S\ref{sec:stm}, where on this
problem it does nothing but sit at a nonzero value.

$H(u)=u^2+\sigma u$ is \emph{strictly} convex ($H_{uu}=2>0$ everywhere) and
$u$ ranges over all of $\R$, so $H_u=0$ is both necessary and sufficient for
the global minimizer: one equation, one formula, the box in \S4. $H(u) =
1+\sigma u$ is \emph{affine} in $u$ ($H_{uu}=0$): an affine function has no
interior minimum, so minimizing it over the compact set $|u|\le u_{\max}$
forces the answer to a boundary point, $u^\star=-u_{\max}\operatorname{sign}\sigma$,
except at the isolated instants where $\sigma=0$ --- a switch. $H(u) =
|u|+\sigma u$ is convex but only \emph{piecewise} linear, with a corner at
$u=0$ and no strict convexity anywhere; minimizing it over $|u|\le u_{\max}$
gives $u^\star=-u_{\max}\operatorname{sign}\sigma$ when $|\sigma|>1$,
$u^\star=0$ when $|\sigma|<1$, and leaves $u^\star$ undetermined by this
condition alone on any interval where $|\sigma|\equiv1$ --- a switch with a
dead zone. In both new problems $\sigma$'s zeros or unit crossings, not a
derivative, decide the structure of the control, and locating them (event
detection, saltation matrices across the jump) is the content of
sub-projects B (minimum time) and C (minimum fuel). Their equations are not
repeated here; this document will get them when the code exists to teach
from, the same rule \S1--\S\ref{sec:stm} followed for this objective.
\end{rig}

\end{document}
